% !TeX root = ../papers/nc-p3.tex
\section{異方向二重色の疎な窓と二変数族の境界}
\label{sec:klein-sparse-transverse}
\begin{theorem}\label{thm:klein-sparse-transverse}
二切断の厚い色の方向が異なる $(2,2,1,1)$ 型の四元群窓を、係数の零を許して考える。
その全正則窓は相異なる二点 $U,V\in\mathbb P^1$ による一つの族で被覆される。
有限な $U,V$ は定理\ref{thm:klein-transverse-dense}の $u,v$ である。
疎な正則窓には、台を保つ対角基底変更と前節の台の作用のもとで、
一つの孤立表示と一つの $y\in k^*$ による表示があり、
前者は $\{U,V\}=\{\infty,0\}$、後者は $\{\infty,1/y\}$ を与える。
全点で最小周期は1で、二点の順序交換と反転は同じ添字保持同型類を与える。
反対階数12の集合は $U=1$, $V=1$, $(U+1)(V+1)=4$ の射影閉包であり、それ以外は16である。
ここでパラメータの全重複を除いたことは仮定しない。
\end{theorem}
\begin{proof}[51の台から二表示へ]
命題\ref{prop:klein-two-doublet-supports}の異方向の51代表ごとに、各非零語に可逆変数を置く。
薄い行列の二次、厚い行列の三次小行列式から、可逆単項式を除いた二項式の指数行列を取る。
51件の非零Smith対角は全て単元である。符号のある右辺との整合を検査し、さらに対角基底変更を割る。
残りの小行列式を代入し、自由変数の積で飽和して極小素分解を取る。
空なものを除いて合計63成分となる。各素成分は、ある変数に関して一次の式を逐次解いて
残る変数を自由にすることで有理化できる。ここで有理化は関数体の計算であり、
上界は元の全縮約列で書いた関係行列の小行列式の消滅として成分全体に延びる。

始点を $0,1,2,3$ と変えて、長さ4の零文字次元を標準値9と比較する。
61成分ではいずれかの始点ですでに10以上となる。
残る台は十六進数 $\mathtt{0ff066663c3caa55}$ と $\mathtt{0ff066667e7eaff5}$ だけである。
第一代表は全非零係数が1、第二代表はSmith自由変数 $y\in k^*$ 一つで追加方程式はない。
両代表の全係数は \path{computations/verify_klein_sparse_transverse.py} の指数変換から再現される。
同検算は全51台の飽和・極小素分解と全63成分の関数体検査を含む。
\end{proof}

\begin{proof}[境界の完全対称表示]
第二代表で $y=t^2$ と書き、四位置の文字順に $t$ の指数
$(1,-1,0,0),(-1,1,0,0),(1,1,0,0),(1,1,0,0)$ を持つ対角倍率を掛ける。
第一代表には倍率を掛けない。どちらも全位置で文字に3を加える作用に不変となる。
従って定理\ref{thm:klein-transverse-dense}と同じ二文字ずつのFourier変換後の総次数は零である。
同じ二コサイクルで捻り、$z=t^{-2}$（第一代表では $z=1$）として、
位置1,2で $(1,z,1,-z),(1,z,1,z)$ を掛け、位置3では文字に1を加えた後に $(1,1,-1,1)$ を掛ける。
全体を2で割ったものは完全対称で、行列 $S_b=(S_{i,i+b,k,k+b+1})$ は
\[
 S_0=[A,A,C,C],\quad S_1=S_0^t,\quad
 S_2=(1,Q,1,Q)^t{\bf1}^t,\qquad S_3=S_2^t,
\]
ただし $A=(1,P,1,Q)^t$, $C=(1,Q,1,P)^t$、
第一代表で $(P,Q)=(3,-1)$、第二代表で $(P,Q)=((3y+1)/y^2,(1-y)/y^2)$ である。

従って平方成分の二条件は $A_0^2=C_0^2$ と $B_0^2+D_0^2=(P+Q)A_0^2$、
他の二式は $Q A_0 C_0=B_0 D_0$, $A_0D_0=B_0 C_0$ である。
$[A_0:B_0:C_0:D_0]=[sm:sn:rm:rn]$ と書く。$m=0$ なら
$s^2+r^2=sr=0$ となって不可能である。$m\ne0$ なら $s^2=r^2\ne0$ となり、
$2n^2=(P+Q)m^2$, $n^2=Qm^2$ から $P=Q$ が必要である。
実際には全点で $P-Q\ne0$ なので、対称二次完全交叉の判定により全て正則である。
\end{proof}

\begin{proof}[射影パラメータと特殊化]
一般に $S_0$ の二種の列を $(a,b,c,d)^t,(c,d,a,b)^t$、
$S_2$ の列を $(c,d,c,d)^t$、$S_1,S_3$ をそれぞれの転置とする。
これが上記の対称表示をまとめた $S$ である。
全位置に同時に $\operatorname{diag}(s,r,s,r)$ を作用させると、
$(a,c)$ は $s^3r$ 倍、$(b,d)$ は $sr^3$ 倍される。従って二対の独立な非零倍率は基底変更で吸収できる。
有限なパラメータを $[a:c]=[U+3:U-1]$, $[b:d]=[V+3:V-1]$ とし、これを射影直線に拡張する。
疎な第一対は $[1:1]$ で $U=\infty$、第二対を $4/(P-Q)$ 倍すれば本文の $V$ が得られる。
この多項式表示が各境界点への特殊化も与える。

元の無印窓は $S$ の係数に四語の二コサイクル符号を掛けたものと四因子同値である。
実際、途中の文字の平行移動を符号の前後で交換して生ずる差は、双指標性により各位置の対角符号の積である。
$S$ の総次数は一定なので、巡回に伴う符号の差も移動させた一文字だけの関数となる。
反転では指数の差が $(\sum i_{\rm hi})(\sum i_{\rm lo})-\sum i_{\rm hi}i_{\rm lo}$ となり、やはり位置ごとに分離する。
従って全点が一周期で反転は同じ類となる。全位置への1の加算は二つの射影パラメータを交換する。
反対階数は有限部分で前定理に従う。疎な曲線では各反対ブロックの行列式の非単元因子は $y-1,y+1$ だけであり、
両零点で各ブロックが階数3、孤立表示では全ブロックが可逆である。これは上の射影閉包に一致する。
\end{proof}
